\section{本章小结}
\label{sec:cdis-summary}

本章开头问：没有概率表、没有密度，或者条件事件的概率为零时，“给定$Y=y$时$X$的分布”是什么。逐事件的条件概率，即指示函数的条件期望，总是存在，它是只依赖观测信息的随机变量，并且几乎必然满足概率的各条规则；但每条规则的例外集依赖所取的事件，不可数多个事件的版本随意挑选，可能在每个观测值处都拼不成分布。

概率核把这件事组织起来：固定观测值，它是一个概率测度；固定事件，它是观测值的可测函数。概率表中它就是行和为一的矩阵。正则条件分布是与联合规律相容的概率核，判别方法是按观测的分布把核积分回去，必须恢复每个矩形事件的联合概率。条件质量函数、条件密度以及$X=Y$时的单点分布$\delta_y$都通过这一判别。

存在性与唯一性都靠把要求压缩到可数个事件上。目标取实数值时，有理阈值上的分布函数值可数，例外集只并可数次，再用下确界补成分布函数；目标取值于标准Borel空间时，用Borel同构搬到实轴上。观测的取值空间不需要任何结构。两个正则条件分布在一个$P_Y$零测集之外相同，但联合规律不决定核在某个零概率观测值上的值。连续性等附加要求可以在零概率点上选出一个值，但满足要求的版本未必存在，选出的值还依赖用哪个观测量逼近，Borel--Kolmogorov悖论中带子与楔形给出$1$与$2x$两种答案。

表~\ref{tab:cdis-summary}对照了条件期望与条件分布。前者对任意可积随机变量存在，回答“给定信息后平均是多少”；后者要求目标空间有标准Borel结构，回答“给定信息后完整的规律是什么”，并通过积分给出前者的一个版本。

\begin{BookInlineTable}
  \centering
  \caption{条件期望与正则条件分布的对照}
  \label{tab:cdis-summary}
  \begin{tabular}{lll}
    \toprule
    & 条件期望$\mathbb{E}[g(X)\mid Y]$ & 正则条件分布$K(y,\cdot)$ \\
    \midrule
    对象 & 一个随机变量 & 每个$y$一个概率测度 \\
    定义 & 已知事件上积分相等 & 每个事件上取条件概率 \\
    存在 & $g(X)$可积即可 & 目标为标准Borel空间 \\
    唯一 & 几乎必然 & 在$P_Y$零测集之外 \\
    关系 & \multicolumn{2}{c}{$\mathbb{E}[g(X)\mid Y]=\int g(x)\,K(Y,\mathrm{d}x)$，$\mathbb{E}\lvert g(X)\rvert<\infty$} \\
    \bottomrule
  \end{tabular}
\end{BookInlineTable}

对核积分有两种方向。固定观测值、对目标积分，得到条件期望的积分表示，证明沿指示函数、简单函数、单调极限与正负分部进行；条件矩与条件方差就是条件分布的矩与方差。整体矩有限只保证条件矩几乎处处有限；反过来，每个观测值上条件矩有限也推不出整体矩有限。再按观测的分布积分，边缘与核唯一地确定联合规律，计算$f(X,Y)$的期望时可以把观测值代入、只对目标积分。对全部观测值积分得到全概率公式，从另一面看就是混合分布：随机硬币的正面次数是二项分布的均匀混合，t分布是高斯分布的尺度混合。

有限个标准Borel随机对象的联合规律总可以链式分解，每一步的核依赖此前的全部变量；反过来，起始分布和这样一列核唯一地确定一个联合规律，对应逐步的随机构造。在链式分解中省去条件变量，需要条件独立作为依据。

条件独立要求给定信息后对所有有界函数的乘积期望分解，等价于条件联合分布是条件边缘的乘积，也等价于已知条件后再知道$X$不改变$Z$的条件分布；条件协方差为零不够。独立与条件独立互不蕴含，增减条件信息都可能改变独立关系；条件化由联合规律决定，干预的效果还依赖产生数据的机制。给定中间变量条件独立时，核可以复合，有限状态下就是转移矩阵相乘；没有条件独立时复合会给出错误的答案。

这些工具在后面多处使用。给定观测后对未知量的分布作修正，是卷八统计推断的出发点，可参考\BookChapterRef{chap:parametric-statistical-inference}{卷八第1章“参数化统计推断”}。把链式分解与核复合沿时间无限地进行下去，就是随机过程中的转移结构，可参考\BookChapterRef{chap:discrete-time-stochastic-processes}{卷九第3章“离散时间随机过程”}。本卷在第~\ref{chap:sample-generation}章按核逐步产生样本，在第~\ref{chap:information-measures}章用条件分布定义条件熵与互信息。
